% !atlas-meta
% id: letter-to-menoeceus
% title: Letter to Menoeceus
% author: epicurus
% authorName: Epicurus
% language: English
% sourceLanguage: Ancient Greek
% translator: Peter Saint-Andre
% translationYear: 2011
% source: https://monadnock.net/epicurus/letter.html
% license: Public Domain (CC0)
% status: complete
% !end-atlas-meta
\documentclass[12pt]{article}
\usepackage[utf8]{inputenc}
\usepackage{ebgaramond}
\usepackage[margin=1.2in]{geometry}
\usepackage{parskip}
\newcommand{\work}[1]{\begin{center}\LARGE\scshape #1\end{center}}
\newcommand{\attrib}[2]{\begin{center}\large #1\\[2pt]\small\itshape trans.\ #2\end{center}\vspace{1em}}
\newcommand{\para}[1][]{\par\noindent\ifx&#1&\else\textsuperscript{\footnotesize #1}~\fi}
\begin{document}
% !body-start
\work{Letter to Menoeceus}
\attrib{Epicurus}{Peter Saint-Andre (public domain)}

\para[1] Greetings from Epicurus to Menoikeus.

\para[2] Let no one put off the love and practice of wisdom when young, nor grow tired of it when old. For it is never too early or too late for the health of the soul. Someone who says that the time to love and practice wisdom has not yet come or has passed is like someone who says that the time for happiness has not yet come or has passed. Young or old, it is necessary to love and practice wisdom, so that in old age you can be youthful by taking joy in the good things you remember, and likewise in youth you can be mature by not fearing what will come. Reflect on what brings happiness, because if you have that you have everything, but if not you will do everything to attain it.

\para[3] Do and practice, then, the things I have always recommended to you, holding them to be the stairway to a beautiful life.

\para[4] First, believe that god is a blissful, immortal being, as is commonly held. Do not ascribe to god anything that is inconsistent with immortality and blissfulness; instead, believe about god everything that can support immortality and blissfulness. For gods there are: our knowledge of them is clear. Yet they are not such as most people believe; indeed most people are not even consistent in what they believe. It is not impious to deny the gods that most people believe in, but to ascribe to the gods what most people believe. The things that most people say about the gods are based on false assumptions, not a firm grasp of the facts, because they say that the greatest goods and the greatest harms come from the gods. For since they are at home with what is best about themselves, they accept that which is similar and consider alien that which is different.

\para[5] Second, train yourself to hold that death is nothing to us, because good and evil consist in sensation, and death is the removal of sensation. A correct understanding that death is nothing to us makes the mortality of life enjoyable --- not because it gives you an unbounded span of time, but because it removes the desire for immortality. There is nothing terrifying in life to someone who truly understands that there is nothing terrifying in the absence of life.

\para[6] Only a fool says that he fears death because it causes pain ahead of time, not because it will cause pain when it comes. For something that causes no trouble when present causes only a groundless pain when merely expected. So death, the most terrifying of evils, is nothing to us, because as long as we exist death is not present, whereas when death is present we do not exist. It is nothing to those who live (since to them it does not exist) and it is nothing to those who have died (since they no longer exist).

\para[7] Most people shrink from death as the greatest of evils, or else extol it as a release from the evils of life. Yet the wise man does not dishonor life (since he is not set against it) and he is not afraid to stop living (since he does not consider that to be a bad thing). Just as he does not choose the greatest amount of food but the most pleasing food, so he savors not the longest time but the span of time that brings the greatest joy. It is simpleminded to advise a young person to live well and an old person to die well, not only because life is so welcome but also because it is through the very same practices that one both lives well and dies well. It is even worse to say that it is good to never have been born, or, having been born, to pass through the gates of Hades as soon as possible.

\para[8] If he believes what he says, why doesn't he depart from life? It is easily done, if he has truly decided. But if he is joking, it is a worthless remark to those who don't accept it. Remember that what will be is not completely within our control nor completely outside our control, so that we will not completely expect it to happen nor be completely disappointed if it does not happen.

\para[9] Third, keep in mind that some desires are natural whereas others are groundless; that among the natural desires some are natural and necessary whereas others are merely natural; and that among the necessary desires some are necessary for happiness, some for physical health, and some for life itself. The steady contemplation of these facts enables you to understand everything that you accept or reject in terms of the health of the body and the serenity of the soul --- since that is the goal of a completely happy life. Our every action is done so that we will not be in pain or fear. As soon as we achieve this, the soul is released from every storm, since an animal has no other need and must seek nothing else to complete the goodness of body and soul. Thus we need pleasure only when we are in pain caused by its absence; but when we are not in pain then we have no need of pleasure.

\para[10] This is why we say that pleasure is the beginning and the end of a completely happy life. For we recognize it as the primary and innate good, we honor it in everything we accept or reject, and we achieve it if we judge every good thing by the standard of how that thing affects us. And because this is the primary and inborn good, we do not choose every pleasure. Instead, we pass up many pleasures when we will gain more of what we need from doing so. And we consider many pains to be better than pleasures, if we experience a greater pleasure for a long time from having endured those pains. So every pleasure is a good thing because its nature is favorable to us, yet not every pleasure is to be chosen --- just as every pain is a bad thing, yet not every pain is always to be shunned. It is proper to make all these decisions through measuring things side by side and looking at both the advantages and disadvantages, for sometimes we treat a good thing as bad and a bad thing as good.

\para[11] Fourth, we hold that self-reliance is a great good --- not so that we will always have only a few things but so that if we do not have much we will rejoice in the few things we have, firmly persuaded that those who need luxury the least enjoy it the most, and that everything natural is easily obtained whereas everything groundless is hard to get. So simple flavors bring just as much pleasure as a fancy diet if all pain from true need has been removed, and bread and water give the highest pleasure when someone in need partakes of them. Training yourself to live simply and without luxury brings you complete health, gives you endless energy to face the necessities of life, better prepares you for the occasional luxury, and makes you fearless no matter your fortune in life.

\para[12] So when we say that pleasure is the goal, we do not mean the pleasure of decadent people and lying in a bed of desire, as is believed by those who are ignorant or who don't understand us or who are ill-disposed to us, but to be free from bodily pain and mental disturbance. For a pleasant life is produced not by drinking and endless parties and bedding boys and women and consuming fish and other delicacies of an extravagant table, but by sober reasoning, searching out the cause of everything we accept or reject, and driving out opinions that cause the greatest trouble in the soul.

\para[13] Practical wisdom is the foundation of all these things and is the greatest good. Thus practical wisdom is more valuable than philosophy and is the source of every other excellence, teaching us that it is not possible to live joyously without also living wisely and beautifully and rightly, nor to live wisely and beautifully and rightly without living joyously. For the excellences grow up together with the pleasant life, and the pleasant life is inseparable from them.

\para[14] In short, whom do you consider better than someone who holds pious opinions about the gods, who is always fearless in the face of death, who has reasoned out the natural goal of life, and who has understood that the limit of good things is easy to fulfill and easy to achieve, whereas the limit of bad things is either short-lived or causes little pain? Someone who laughs at destiny, which is asserted by some to be the master of all things? For he holds that we are responsible for what we achieve, even though some things happen by necessity, some by chance, and some by our own power, because although necessity is not accountable he sees that chance is unstable whereas the things that are within our power have no other master, so that naturally praise and blame are inseparably connected to them. Indeed he sees that it would be better even to cleave to the myths about the gods (since that leaves some hope of prevailing upon them through worship) than to be subject to the destiny of the scientists (since that way lies an inexorable necessity). And such a man holds that Fate is not a god (as most people believe) because a god does nothing disorderly, and he holds that Fate is not an uncertain cause because nothing good or bad with respect to a completely happy life is given to men by chance, although it does provide the beginnings of both great goods and great evils. And he considers it better to be rationally unfortunate than irrationally fortunate, since it is better for a beautiful choice to have the wrong results than for an ugly choice to have the right results just by chance.

\para[15] So practice these and similar things day and night, by yourself and with a like-minded friend, and you will never be disturbed whether waking or sleeping, and you will live as a god among men: for a man who lives in the midst of immortal goods is unlike a merely mortal being.
% !body-end
\end{document}
